%<<CONT
% !PRINTED 228
\weberpage{228}
is, however, if we set
\[
\Phi =
\begin{vmatrix}
1, & F_0(x_1), & F_1(x_1) & \ldots & F_{n-2}(x_1) \\
1, & F_0(x_2), & F_1(x_2) & \ldots & F_{n-2}(x_2) \\
\multicolumn{5}{c}{\cdots\cdots\cdots\cdots\cdots\cdots\cdots\cdots} \\
1, & F_0(x_n), & F_1(x_n) & \ldots & F_{n-2}(x_n)
\end{vmatrix}
\]
by (11)
\weq{12}{\Phi^2 = n \varDelta.}
If one now takes account of the expressions (9) in \S.\,68 for the functions
$F_0, F_1 \ldots F_{n-2}$, then by the theorem that in a determinant
a column multiplied by an arbitrary factor may be added to
another column (\S.\,22, VII), one easily obtains for $\Phi$ the expression
\[
a_0^{n-1}
\begin{vmatrix}
1, & x_1, & x_1^2 & \ldots & x_1^{n-1} \\
1, & x_2, & x_2^2 & \ldots & x_2^{n-1} \\
\multicolumn{5}{c}{\cdots\cdots\cdots\cdots\cdots\cdots} \\
1, & x_n, & x_n^2 & \ldots & x_n^{n-1}
\end{vmatrix},
\]
whose square, by \S.\,46, is the discriminant of $f(x)$. Accordingly,
by (12), we thus have the important theorem:

\gesperrt{The determinant of the Bezoutiant of $f(x)$ is the
$n$th part of the discriminant of $f(x)$.}

The computation of the Bezoutiant presents, by our formulas,
no further difficulty at all, and turns out to be fairly simple.
For $n = 3$ one obtains the result, already to be inferred from
the formulas of \S.\,71,
\weq{13}{B = - \frac{2}{3} H(t_1, t_0).}
For the computation of the Bezoutiant of the biquadratic
form we shall, to illustrate the foregoing, set out the
system of formulas in full, leaving the carrying out of the
slightly longer computation for the form of the fifth degree
to the reader.

For the form of the fourth degree we have, by (7) and \S.\,72, (9)
\[
f(x) = a_0 x^4 + a_1 x^3 + a_2 x^2 + a_3 x + a_4
\]
\[
\begin{array}{rcl|r}
E_{0,1} &=& -a_2 t_2 - a_3 t_1 - a_4 t_0 & 4 a_0 \tau_2 \\
E_{1,1} &=& a_1 t_2 & 3 a_1 \tau_2 \\
E_{2,1} &=& a_0 t_2 + a_1 t_1 & 2 a_2 \tau_2 \\
E_{3,1} &=& a_0 t_1 + a_1 t_0 & a_3 \tau_2 \\
\hline
\end{array}
\]
%CONT>>
